% Source: mise.tex
\begin{notechapter}
% ---- chapter-local definitions ----
\tikzset{%
    myarrow/.style = {-Stealth, shorten >=5pt}
}
\forgetcs\mypoint\newcommand{\mypoint}[2]{\tikz[remember picture]{\node[inner sep=0, anchor=base](#1){$#2$};}}
\forgetcs\framedtext\newcommand{\framedtext}[1]{%
\par%
\noindent\fbox{%
    \parbox{\dimexpr\linewidth-2\fboxsep-2\fboxrule}{#1}%
}%
}
\forgetcs\horzbar\newcommand*{\horzbar}{\rule[.5ex]{5.5ex}{0.5pt}}
\forgetcs\vertbar\newcommand*{\vertbar}{\rule[-1ex]{0.5pt}{5.5ex}}
% -----------------------------------
\chapter{Mean Integrated Squared Error}\label{ch:mise}

    The Mean Integrated Squared Error (MISE) of a density estimator $\hat{f}$ of an unknown density $f$ is given by:
    \begin{equation}
    MISE(\hat{f})= \mathbb{E}_\mathbf{\mathcal{T}} \left[\int \left(\hat{f}_{\mathcal{T}}(\mathbf{x}) - f(\mathbf{x})\right)^2 d\mathbf{x}\right]
    \end{equation}
    Where the subscript $\mathbf{\mathcal{T}}$ is a reminder that the expectation is over the sampled data.\\[6pt]
    That is $\mathbf{\mathcal{T}} = \{\mathbf{X}_{1}, \dots, \mathbf{X}_{n}\}$ where 
    $\mathbf{X}_{1}, \dots, \mathbf{X}_{n}$ are i.i.d random variables with $\mathbf{X}_i \thicksim f$. \\[6pt]
    %
    Accordingly we will note $\tau$ a realization of $\mathcal{T}$.\\[6pt]
    $\tau =\{\mathbf{x}_1, \mathbf{x}_2, \dots, \mathbf{x}_n\}$ with 
    $\mathbf{x}_i$ a realization of $\mathbf{X}_i$.
    \\[10pt]
    Minimizing $MISE(\hat{f})$ is equivalent to minimizing 
    $\mathbb{E}_\mathbf{\mathcal{T}} \bigg[ \displaystyle\int \hat{f}_{\mathcal{T}}^{\, 2}(\mathbf{x})d\mathbf{x} -2 \displaystyle\int \hat{f_{\mathcal{T}}}(\mathbf{x}) f(\mathbf{x}) d\mathbf{x} \bigg]$.
    \\[10pt]
    %%
    We want to compute the expectations over all training set
    $\mathbb{E}_\mathbf{\mathcal{T}} \bigg[ \displaystyle\int \hat{f}_{\mathcal{T}}^{\, 2}(\mathbf{x})d\mathbf{x} \bigg]$
    and
    $\mathbb{E}_\mathbf{\mathcal{T}} \bigg[ \displaystyle\int \hat{f_{\mathcal{T}}}(\mathbf{x}) f(\mathbf{x}) d\mathbf{x} \bigg]$.
    \\[10pt]
    Since we do not have access to $f$, we wish to compute estimates of these two expectations.\\[10pt]
    A straightforward unbiased estimator of 
    $\mathbb{E}_\mathbf{\mathcal{T}} \bigg[ \displaystyle\int \hat{f}_{\mathcal{T}}^{\, 2}(\mathbf{x})d\mathbf{x} \bigg]$
     is simply 
    $\displaystyle\int \hat{f}_{\mathcal{T}}^{\, 2}(\mathbf{x})d\mathbf{x}$.\\[6pt]
    Therefore, $\displaystyle\int \hat{f}_{\tau}^{\, 2}(\mathbf{x})d\mathbf{x}$ is going to be our estimate of 
    $\mathbb{E}_\mathbf{\mathcal{T}} \bigg[ \displaystyle\int \hat{f}_{\mathcal{T}}^{\, 2}(\mathbf{x})d\mathbf{x} \bigg]$.\\[10pt]
    For $\mathbb{E}_\mathbf{\mathcal{T}} \bigg[ \displaystyle\int \hat{f_{\mathcal{T}}}(\mathbf{x}) f(\mathbf{x}) d\mathbf{x} \bigg]$
    things are a bit more complex.\\[6pt]
    First we notice that we have:\\[6pt]
    $\displaystyle\int \hat{f}_{\tau}(\mathbf{x}) f(\mathbf{x}) d\mathbf{x} = 
    \mathbb{E}_\mathbf{X} \Big[ \hat{f}_{\tau}(\mathbf{X}) \Big]
    $ with $\mathbf{X} \thicksim f$\\[10pt]
    Hence, 
    \begin{flalign*}
    \mathbb{E}_{\mathbf{\mathcal{T}}} \bigg[ \displaystyle\int \hat{f}_{\mathcal{T}}(\mathbf{x}) f(\mathbf{x}) d\mathbf{x} \bigg] &=
    \mathbb{E}_{\mathbf{\mathcal{T}}} \bigg[\mathbb{E}_\mathbf{X} \Big[ \hat{f}_{\mathcal{T}}(\mathbf{X})\, \big| \,\mathcal{T} = \tau\Big]\bigg]
    \end{flalign*}
    \\
    In what follows $\hat{f}$ is a kernel density estimator, $\hat{f}_{\tau}(\mathbf{x}) = \dfrac{1}{n}\displaystyle\sum_{i=1}^{n} \Phi_{h}(\mathbf{x}, \mathbf{x}_{i})$, where $\Phi_h$ is a kernel of bandwidth $h$.
    We are going to use \textbf{leave-one-out cross-validation} to compute our estimate.
    \\[6pt]
    Let's split $\tau$ into $n$ folds of size 1. $\tau = \bigcup\limits_{i=1}^{n} \{\mathbf{x}_{i}\}$.
    And we define $\tau_{\text{-}k} = \bigcup\limits_{\substack{i=1 \\ i\neq k}}^{n} \{\mathbf{x}_{i}\}$.\\[6pt]
    %%
    Call $\hat{f}_{\text{-}k}$ the estimate we get using $\tau_{\text{-}k}$. i.e. 
    $\hat{f}_{\text{-}k}(\mathbf{x}) = \dfrac{1}{n-1} \displaystyle\sum\limits_{\substack{i=1 \\ i\neq k}}^{n} \displaystyle{\Phi}_{h}(\mathbf{x}, \mathbf{x}_{i})$\\[10pt]
    %%
    Then, $\dfrac{1}{n} \displaystyle\sum\limits_{k=1}^{n} \hat{f}_{\text{-}k}(\mathbf{x}_{k})$ 
    is an unbiased estimate of 
    $\mathbb{E}_{\mathbf{\mathcal{T}}} \bigg[\mathbb{E}_\mathbf{X} \Big[ \hat{f}_{\mathcal{T}}(\mathbf{X})\, \big| \,\mathcal{T} = \tau\Big]\bigg]$.
    \\[10pt]
    \paragraph{Proof 1.}

    \begin{flalign*}
        \begin{WithArrows}
        \mathbb{E} \bigg[  \dfrac{1}{n} \displaystyle\sum\limits_{k=1}^{n} \hat{f}_{\text{-}k}(\mathbf{X}_{k}) \bigg] &=
        \dfrac{1}{n} \displaystyle\sum\limits_{k=1}^{n} \mathbb{E} \big[ \hat{f}_{\text{-}k}(\mathbf{X}_{k}) \big] \\
        &=\dfrac{1}{n} \displaystyle\sum\limits_{k=1}^{n} \mathbb{E}_\mathbf{\mathcal{T_{\text{-k}}}} 
        \bigg[ \mathbb{E}_\mathbf{\mathcal{\mathbf{X}}_{k}} \Big[\hat{f}_{\text{-}k}(\mathbf{X}_{k}) \, \Big| \,\mathcal{T}_{\text{-}k} = 
        \tau_{\text{-}k} \Big] \bigg] \Arrow[tikz={->, in=100, out=80, looseness=1.5}]{\text{$\hat{f}$ is an average of i.i.d.\ terms: its expectation does not depend on the sample size}} \\
        &=\dfrac{1}{n} \displaystyle\sum\limits_{k=1}^{n} \mathbb{E}_\mathbf{\mathcal{T}} \bigg[ \mathbb{E}_\mathbf{\mathcal{\mathbf{X}}} \Big[\hat{f}(\mathbf{X})  \, \Big| \,\mathcal{T} = \tau \Big] \bigg] \\
        &=\mathbb{E}_{\mathbf{\mathcal{T}}} \bigg[\mathbb{E}_\mathbf{X} \Big[ \hat{f}_{\mathcal{T}}(\mathbf{X})\, \Big| \,\mathcal{T} = \tau\Big]\bigg]
        \end{WithArrows}
    \end{flalign*}\\
    conclusion: \\[10pt]
    $\dfrac{1}{n} \displaystyle\sum\limits_{k=1}^{n} \hat{f}_{\text{-}k}(\mathbf{X}_{k})$ is an unbiased estimator of
    $\mathbb{E}_\mathbf{\mathcal{T}} \bigg[ \displaystyle\int \hat{f_{\mathcal{T}}}(\mathbf{x}) f(\mathbf{x}) d\mathbf{x} \bigg]$,
    and $\dfrac{1}{n} \displaystyle\sum\limits_{k=1}^{n} \hat{f}_{\text{-}k}(\mathbf{x}_{k})$ is our estimate.
    %
    \\[15pt]
    \paragraph{Proof 2 (using the expression of the kernel estimator).}

    \begin{flalign*}
        \mathbb{E} \bigg[  \dfrac{1}{n} \displaystyle\sum\limits_{k=1}^{n} \hat{f}_{\text{-}k}(\mathbf{X}_{k}) \bigg] &=
        \dfrac{1}{n} \displaystyle\sum\limits_{k=1}^{n} \mathbb{E} \big[ \hat{f}_{\text{-}k}(\mathbf{X}_{k}) \big] \\
        &= \dfrac{1}{n} \displaystyle\sum\limits_{k=1}^{n} \bigg[  
            \dfrac{1}{n-1} \displaystyle\sum\limits_{\substack{i=1 \\ i\neq k}}^{n} \mathbb{E} \Big[ 
                \displaystyle{\Phi}_{h}(\mathbf{X}_{k}, \mathbf{X}_{i}) \Big] \bigg] \\
        &= \dfrac{1}{n} \displaystyle\sum\limits_{k=1}^{n} \bigg[ 
            \dfrac{1}{n-1} \displaystyle\sum\limits_{\substack{i=1 \\ i\neq k}}^{n} \Big[ 
                \displaystyle\int  \displaystyle\int \displaystyle{\Phi}_{h}(\mathbf{x}, \mathbf{y}) f(\mathbf{x}) f(\mathbf{y}) d\mathbf{x} d\mathbf{y} \Big] \bigg]\\
        &= \displaystyle\int  \displaystyle\int \displaystyle{\Phi}_{h}(\mathbf{x}, \mathbf{y}) f(\mathbf{x}) f(\mathbf{y}) d\mathbf{x} d\mathbf{y} \\
        &= \displaystyle\int  \bigg[ \displaystyle\int \displaystyle{\Phi}_{h}(\mathbf{x}, \mathbf{y}) f(\mathbf{y}) d\mathbf{y} \bigg] f(\mathbf{x}) d\mathbf{x} \\
        &= \displaystyle\int  \mathbb{E} \big[  \hat{f}_{\mathbf{x}} (\mathbf{Y}) \big] f(\mathbf{x}) d\mathbf{x} \\
        &= \mathbb{E} \bigg[ \displaystyle\int \hat{f}(\mathbf{x}) f(\mathbf{x}) d\mathbf{x} \bigg]
    \end{flalign*}
\end{notechapter}
